\subsection{极分解}

引理 8. --- 设 E 是分离的拓扑向量空间。设 \((\mathrm{E}_{1},\|\cdot\|_{1})\) 与 \((\mathrm{E}_{2},\|\cdot\|_{2})\) 是 E 的稠密子空间，其上赋以希尔伯特结构，使得 \(E_{1}\) 与 \(E_{2}\) 到 E 中的典范单射都是连续的。设 F 是 \(E_{1}\cap E_{2}\) 的子空间，它对于这些希尔伯特结构在 \(E_{1}\) 与 \(E_{2}\) 中稠密。若 \(\|x\|_{1}=\|x\|_{2}\) 对所有 \(x\in F\) 成立，则 \(E_{1}=E_{2}\) 且 \(\|\cdot\|_{1}=\|\cdot\|_{2}\) 。

设 \(x \in E_{1}\) 。存在 F 中的序列 \((x_{n})_{n \in \mathbf{N}}\) 使得 \(x_{n}\) 在希尔伯特空间 \(E_{1}\) 中收敛于 x。由于对所有整数 n 与 m 有 \(\|x_{n} - x_{m}\|_{2} = \|x_{n} - x_{m}\|_{1}\) ，序列 \((x_{n})\) 是 \(E_{2}\) 中的柯西序列。用 y 表示它的极限。我们有 \(\|y\|_{2} = \lim\|x_{n}\|_{2} = \|x\|_{1}\) 。由于 \(E_{1}\) 与 \(E_{2}\) 到 E 中的典范单射都是连续的，我们有 \(x_{n} \to x\) 在 E 中成立，同样 \(x_{n} \to y\) 也在 E 中成立。于是由此得出 x = y 且 \(\|x\|_{1} = \|x\|_{2}\) 。特别地，\(E_{1} \subset E_{2}\) 。由对称性即得结论。

设 E 是一个复希尔伯特空间。设 u 是 E 上的正自伴部分算子。对每个 \(\alpha \in \mathbf R_{+}\) ，我们置 \(u^{\alpha} = f(u)\) ，其中 f 是连续映射 \(x \mapsto x^{\alpha}\) ，从 \(\mathbf R_{+}\) 到 \(\mathbf R\) 中。这是一个正自伴部分算子（IV, p. 273 的推论 1, a)）。当 \(u \in \mathcal{L}(E)\) 时，这个记号与关于 C*-代数 \(\mathcal{L}(E)\) 的记号相容（参见 I, p. 118 的命题 16）。若 \(\alpha\) 是正整数，则部分算子 \(u^{\alpha}\) 与由复合 \(u \circ \cdots \circ u\) 所定义的部分算子重合（IV, p. 273 的推论 1, b)）。

设 \(\beta \in \mathbf{R}_{+}\) 。我们有 \(u^{\alpha \beta} = (u^{\alpha})^{\beta}\) （IV, p. 274 的推论 4）。特别地，若 \(\alpha > 0\) ，则部分算子 \(u^{1 / \alpha}\) 是 E 上唯一的正自伴部分算子 \(v\) ，使得 \(v^{\alpha} = u\) 。

假设 \(0 \leqslant \alpha \leqslant \beta\) 。这时我们有 \(\operatorname{dom}(u^{\beta}) \subset \operatorname{dom}(u^{\alpha})\) ：事实上，对每个实数 \(x \geqslant 0\) ，有 \(x^{\alpha} \leqslant 1 + x^{\beta}\) ，于是结论由 \(u^{\alpha}\) 的定义域的定义得出（参见 IV, p. 271 的命题 3）。

设 u 是从 E 到复希尔伯特空间 F 中的闭稠定算子。部分算子 \(u^{*}u\) 是自伴且正的（IV, p. 241 的命题 12）。我们记 \(|u| = (u^{*}u)^{1/2}\) 。

命题 12. --- 设 u 是从 E 到复希尔伯特空间 F 中的闭稠定算子。

\begin{enumerate}
\def\labelenumi{\alph{enumi})}
\item
  正自伴部分算子 \(|u|\) 的定义域与 \(u\) 的定义域重合；
\item
  存在唯一的从 E 到 F 中的部分等距线性映射 \(j\) ，使得 \(u = j|u|\) 且 \(\operatorname{Ker}(j) = \operatorname{Ker}(u)\) ；
\item
  设 \(u_{1}\) 是 E 上的正自伴算子，\(j_{1}\) 是从 E 到 F 中的部分等距线性映射，使得 \(u = j_{1}u_{1}\) 且 \(\operatorname{Ker}(j_{1}) = \operatorname{Ker}(u_{1})\) 。这时我们有 \(u_{1} = |u|\) 且 \(j_{1} = j\) 。
\end{enumerate}

\(u^{*}u\) 的定义域包含在 \(\operatorname{dom}(u)\) 与 \(\operatorname{dom}(|u|)\) 中。它在希尔伯特空间 \(E_{u}\) （loc. cit.）与 \(E_{|u|}\) （IV, p. 273 的推论 1, e)）中稠密，且对每个 \(x \in \operatorname{dom}(u^{*}u)\) ，我们有

\[
\begin{array}{c} \| x \| _ {u} ^ {2} = \| x \| ^ {2} + \langle u (x) | u (x) \rangle = \| x \| ^ {2} + \langle x | (u ^ {*} u) (x) \rangle \\ = \| x \| ^ {2} + \langle | u | (x) | | u | (x) \rangle = \| x \| _ {| u |} ^ {2}. \end{array}\tag{16}
\]

因此，希尔伯特空间 \(E_{u}\) 与 \(E_{|u|}\) 相等（引理 8），从而 \(\text{dom}(u) = \text{dom}(|u|)\) ，且 \(\|u(x)\| = \||u|(x)\|\) 对所有 \(x \in \text{dom}(u)\) 成立。

公式 (16) 蕴含 \(\operatorname{Ker}(u) = \operatorname{Ker}(|u|)\) ，并且存在唯一的等距线性映射 \(v\) ，它从 \(\overline{\operatorname{Im}(|u|)}\) 到 \(\overline{\operatorname{Im}(u)}\) 上，且满足 \(v(|u|(x)) = u(x)\) 对所有 \(x \in \operatorname{dom}(|u|)\) 成立。由于 \(|u|\) 是自伴的，我们有 \(\operatorname{Im}(|u|)^{\circ} = \operatorname{Ker}(|u|)\) （IV, p. 236 的命题 7, c)）。因此存在唯一的从 E 到 F 中的部分等距 \(j\) ，它延拓 \(v\) 且在 \(\operatorname{Ker}(|u|) = \operatorname{Ker}(u)\) 上为零。由于 \(\operatorname{E} = \operatorname{Ker}(u) \oplus \overline{\operatorname{Im}(|u|)}\) ，这个映射是使 \(u = j|u|\) 且 \(\operatorname{Ker}(j) = \operatorname{Ker}(u)\) 的唯一的部分等距映射。

我们来证明 c)。我们有 \(u_{1}j_{1}^{*}j_{1}u_{1}\subset(j_{1}u_{1})^{*}j_{1}u_{1}=u^{*}u\) 。线性映射 \(j_{1}^{*}j_{1}\) 是以 \(\mathrm{Ker}(j_{1})=\mathrm{Ker}(u_{1})\) 为核（EVT, V, p. 41, 命题 5 (ii)），从而以 \(\mathrm{Ker}(u_{1})^{\circ}=\overline{\mathrm{Im}(u_{1})}\) 为像（IV, p. 236 的命题 7, c)）的正交投影算子。因此，\(u_{1}^{2}\subset u^{*}u\) ，又由于这两个算子都是自伴的，故 \(u_{1}^{2}=u^{*}u\) 。于是由此得出 \(u_{1}=(u^{*}u)^{1/2}\) ，而 b) 中的唯一性断言最终表明 \(j_{1}=j\) 。

定义 7. --- 设 u 是从 E 到复希尔伯特空间 F 中的稠定闭算子。由命题 12 所确定的偶 \((j, |u|)\) 称为 u 的极分解。

假设 \(u \in \mathcal{L}(\mathrm{E};\mathrm{F})\) 。它在这个定义意义下的极分解与 I, p. 140 的定义 4 中的极分解重合。

